\documentclass[utf8,bachelor]{gradu3}

\usepackage{graphicx} % kuvien mukaan ottamista varten

\usepackage{amsmath} % hyödyllinen jos tekstisi sisältää matikkaa,
                     % ei pakollinen

\usepackage{booktabs} % hyvä kauniiden taulukoiden tekemiseen

% HUOM! Tämän tulee olla viimeinen \usepackage koko dokumentissa!
\usepackage[bookmarksopen,bookmarksnumbered,linktocpage]{hyperref}

\addbibresource{veikkakumpulainen_kandi.bib} % Lähdetietokannan tiedostonimi

\begin{document}

\title{Tutkimussuunnitelma}
\translatedtitle{Research Plan}
\studyline{Tieto- ja ohjelmistotekniikka}
\avainsanat{%
  \LaTeX,
  kandidaatintutkielmat}
\keywords{\LaTeX, Bachelor's Theses}
\tiivistelma{%
Tutkimus generatiivisen tekoälyn vaaroista ja huonoista puolista ohjelmistokehityksessä.
}
\abstract{%
Research on the risks and disadvantages of using generative artificial intelligence in software development
}

\author{Veikka Kumpulainen}
\contactinformation{\texttt{veikka.e.kumpulainen@jyu.fi}}
\supervisor{Sanna Juutinen}
\subject{Tieto- ja ohjelmistotekniikan}

\maketitle

Jyväskylässä \today

\bigskip

Tutkielman tekijä

%\begin{thetermlist}
%\item[\TeX] Donald Knuthin 1977--1989 laatima eräajotyyppinen
% ladontajärjestelmä \parencite[ks.][]{knuth86:_texbook}. 
%\item[\LaTeX] \TeX in \parencite[ks.][]{knuth86:_texbook} päälle
% rakennettu rakenteisten kirjoitelmien ladontaan tarkoitettu
%  järjestelmä \parencite[ks.][]{lamport94:_latex}.  Siitä on nykyään
%  käytössä versio \LaTeXe.
%\end{thetermlist}

\mainmatter

\chapter{Johdanto}

Generatiivista tekoälyä käytetään nykyään paljon ohjelmistokehityksessä, mutta tietääkö kaikki ohjelmistokehityksen parissa työskentelevät sen vaaroista ja huonoista puolista. Teen kirjallisuuskatsauksen siitä, mitä vaaroja ja huonoja puolia generatiivisen tekoälyn käytöstä voi seurata ohjelmistokehityksessä. Tutkin, koska kaikkien ohjelmistokehityksen parissa työskentelevien olisi hyvä tietää haitoista. Tutkin lukemalla aiheesta tehtyjä tutkimuksia ja muodostamalla niiden pohjalta kokonaisuuden vaaroista ja huonoista puolista.

\chapter{Kirjallisuuskartoitus}

Ohelmistokehityksellä tarkoitetaan joukkoa toimintoja, jotka keskittyvät ohjelmiston luomiseen, suunniteluun, käyttöönottoon ja tukemiseen. Ohjelmistokehityksen tavoitteena on luoda tuote, joka täyttää käyttäjien tarpeet ja liiketoiminnan tehokkaalla, toistettavalla ja turvallisella tavalla. Ohjelmistokehittäjät, ohjelmoijat ja ohjelmistosuunnittelijat kehittävät ohjelmistoja sarjassa, jota kutsutaan ohjelmistokehityksen elinkaareksi. Tekoälypohjaisia työkaluja ja generatiivista tekoälyä käytetään yhä enemmän tukemaaan ohjelmistokehitystiimejä koodin tuottamisessa ja testaamisessa. \parencite{noauthor_what_2024}


Generatiivinen tekoäly tarkoittaa syväoppimismalleja, jotka pystyvät tuottamaan monimutkaistaa alkuperäistä sisältöä, kuten pitkiä tekstikokonaisuuksia, korkealaatuisia kuvia, realistisia videoita tai ääntä sekä muuta sisältöä käyttäjän antaman kehotteen perusteella. Yleisellä tasolla generatiiviset mallit koodavat yksinkertaistetun esityksen koulutusdatastaan ja hyödyntävät tätä esitystä tuottaakseen uutta sisältöä, joka on samankaltaisia mutta ei identtisitä alkuperäisen datan kanssa. \parencite{noauthor_what_2024-1} Eli generatiivinen tekoäly pystyy luomaan myös koodia.

Kun tekoälyä käyttää ohjelmistonkehitykseen, se luo piilevää teknistä velkaa. Vaikka ongelma ratkeaa nopeasti ei tiedetä, mitä ongelmia nopea tekoälyn suosittelema korjaus tuo ja mahdollisia kuluja jälkikäteen. \parencite{sculley_hidden_2015} 

GutHub Copilotin generoima koodi tuottaa toimiviaa koodia, mutta usein häviää suorituskyvyssä ihmisen tekemälle. \parencite{li_assessing_2024}

\chapter{Tutkimusaihe ja tutkimuskysymykset}

Tutkimusaihe on generatiivisen tekoälyn käytön riskit ja huonot puolet ohjelmistokehityksessä, joista ohjelmistonkehityksen parissa työskentelevien olisi hyvä olla tietoinen. Tutkimuskysymyksiä voisi olla mitä haittoja generatiivisen tekoälyn vastuuton käyttö aiheuttaa ohjelmistokehityksessä. Mitä vaaroja liittyy tekoälyllä generoidun koodin luotettavuuteen ja turvallisuuteen. Miten tekoälyn haitat voidaan minimoida ohjelmistokehityksessä.

\chapter{Tutkimusmenetelmä}

Etsin lähteitä scopuksesta eri hakusanoilla ja lisäksi kysyin copilotilta hyviä mahdollisia lähteitä koskien aihettani. Googlesta etsin suoraan myös lähteitä. Valitsin parhaat lähteet sen perusteella, jos niissä oli käsitelty aiheeni kannalta ajankohtaisia asioita ja julkaisualusta oli luotettava. Lisäksi löysin yhden kandityön samaan aiheeseen liittyen ja katsoin sen lähdeluettelosta mahdollisia hyviä lähteitä.

\chapter{Johtopäätökset ja tutkimuksen merkitys}

Tutkimuksen avulla ohjelmistokehityksen parissa työskentelevät voivat saada paremman tietoisuuden riskeistä ja huonoista puolista, joita generatiivisen tekoälyn vastuuttomasta käytöstä voi seurata ja parantaa omaa työnlaatua ja koodinlaatua tämän avulla.

\printbibliography

\end{document}
